\noindent\hspace*{\parindent}Dem ist nun der Kapitalbetrag der Einnahmequellen der Witwenkasse gegenüberzustellen. Da die Beiträge der Mitglieder beim ersten obigen Ansatz bereits abgezogen waren, sind hier nur noch das Geldvermögen und die Apothekenpacht zu veranschlagen.

Das Vermögen ist im letzten Rechnungsabschluss, also für 1845 Juli 1, mit 116369 Thl. 17 Ggr. 4 Pf. ausgewiesen. Um es jedoch auf den 1. Oktober entsprechend der oben S. [149] angeführten Rechnungsgrundlage umzurechnen, muss der zu Michaelis 1845 fällig gewordene halbjährliche Betrag der Witwenpensionen von 2350 Thl. abgezogen werden. Hinzuzufügen sind dagegen die Beiträge der Mitglieder für das Jahr 1844—1845, nämlich volle Beiträge von 480 Thl. und Teilzahlungen von 18 Thl. 17 Ggr. 4 Pf., ferner die zu Michaelis 1845 fällig gewordene halbjährliche Apothekenpacht von 500 Thl. Um nichts zu übergehen, müssten auch noch die während der drei Monate vom 1. Juli bis zum 1. Oktober eingegangenen Zinsen hinzugerechnet werden, deren Betrag mir unbekannt ist. Es wird aber wohl nicht viel fehlen, wenn ich dafür den vierten Teil der Zinseinnahmen des letzten Jahres nach Abzug der Einziehungskosten mit 1163 Talern in Rechnung stelle. Danach setze ich das Vermögen der Kasse am 1. Oktober 1845 unter Weglassung der Bruchteile mit 116171 Talern an.

Die Apothekenpacht in ihrer jetzigen Höhe von 1000 Thl. entspricht einem Kapitalvermögen von 28571 Thl. oder 25000 Thl., je nach der Höhe des vorausgesetzten Zinsfußes. Das effektive Vermögen der Kasse beträgt folglich:
\begin{center}
144742 Thl. oder 141171 Thl., je nachdem man \(3\tfrac{1}{2}\) oder 4 Prozent annimmt.
\end{center}
\noindent\hspace*{\parindent}Daraus ergibt sich schließlich die \textit{Bilanz} der Kasse
\begin{center}
\begin{tabular}{l}
als ein \textit{Defizit} von 13990 Thl. für Zinsfuß \(3\tfrac{1}{2}\) Prozent\\
als ein \textit{Überschuss} von 2508 Thl. für Zinsfuß 4 Prozent
\end{tabular}
\end{center}
\noindent\hspace*{\parindent}Dass der unter der zweiten Voraussetzung entstehende Überschuss bei Weitem kleiner ist als die bei den einzelnen Bestandteilen der Rechnung verbleibenden Unsicherheiten, braucht wohl kaum bemerkt zu werden: Ein einziger Todesfall kann die ganze Bilanz leicht auf einmal um 3000—3500 Rthl. zum Nachteil der Kasse verändern. Außerdem darf aber auch nicht übersehen werden, dass ich für mögliche Verluste durch Zahlungsunfähigkeit der Schuldner überhaupt nichts angesetzt habe: Ich habe dies unterlassen, weil jede Schätzung im Voraus äußerst unsicher bleiben muss. Es wäre interessant, alle seit der Gründung eingetretenen Verluste zusammenzustellen; bei der Form der Jahresrechnungen, zumal in früheren Zeiten, wäre dies aber nicht ohne einen sehr großen Zeitaufwand ausführbar. Ich selbst habe mich, um doch einigermaßen eine Vorstellung davon zu erhalten, mit einer Zusammenstellung der letzten 14 Jahre begnügen müssen. Sie hat (falls ich nichts übersehen habe) einen Gesamtverlust von 1729 Thl. 14 Ggr. 1 Pf., also einen durchschnittlichen jährlichen Verlust \(=123\) Thl. 13 Ggr. ergeben. Kapitalisiert beträgt dies je nach dem Zinsfuß 3530 Thl. oder 3088 Thl. Dürfte man diesen Durchschnitt als maßgebend betrachten, so würde sich die \textit{Bilanz} für \textit{beide} Zinsfüße als ein \textit{Defizit} ergeben, nämlich
\begin{center}
\begin{tabular}{l}
von 17520 Thl. bei dem Zinsfuß von \(3\tfrac{1}{2}\) Prozent\\
von\phantom{ 17} 580 Thl. bei dem Zinsfuß von 4 Prozent.
\end{tabular}
\end{center}
\noindent Göttingen 1845 November 2.\hfill C. F. \textsc{Gauss}.

